\documentclass[11pt]{article}
\usepackage[a4paper,margin=1in]{geometry}
\usepackage{enumitem}
\usepackage{hyperref}
\usepackage{amsmath}
\usepackage{booktabs}
\usepackage{longtable}
\usepackage{array}

\title{Mental Math --- Level Definitions (v1)}
\date{}

\begin{document}

\maketitle

\section{Overview}

Each domain (Addition, Subtraction, Multiplication, Division) is divided into ten progressive levels.
A \textbf{pattern} defines the shape of questions generated for that level: the number of digits in each
operand, the required number of carries or borrows, and the cutoff time within which a correct answer
must be submitted to count as fast.

Levels are designed to be completed in order within a domain. Progression to the next level requires
demonstrating consistent performance on the current one (see \texttt{progression-algo/algorithm.tex}).

\bigskip
\noindent Column definitions used in the tables below:

\begin{itemize}[leftmargin=*]
  \item \textbf{Level} --- sequential difficulty step within the domain (1 = easiest).
  \item \textbf{Description} --- plain-language summary of the pattern.
  \item \textbf{Operands} --- digit widths of the left and right operands (or dividend / divisor for division).
  \item \textbf{Carries / Borrows} --- exact number of carry or borrow operations required by the generated question.
  \item \textbf{Cutoff (ms)} --- time budget per question; answers submitted within this window are considered fast.
\end{itemize}

% ──────────────────────────────────────────────
\section{Addition (ADD)}

\begin{longtable}{@{}clccr@{}}
\toprule
\textbf{Level} & \textbf{Description} & \textbf{Operands} & \textbf{Carries} & \textbf{Cutoff (ms)} \\
\midrule
\endhead
1  & 1-digit + 1-digit, no carry          & $1 + 1$ & 0 & 4\,000 \\
2  & 1-digit + 1-digit, with carry        & $1 + 1$ & 1 & 4\,000 \\
3  & 1-digit + 2-digit, no carry          & $1 + 2$ & 0 & 4\,500 \\
4  & 1-digit + 2-digit, with carry        & $1 + 2$ & 1 & 4\,500 \\
5  & 2-digit + 2-digit, no carry          & $2 + 2$ & 0 & 5\,000 \\
6  & 2-digit + 2-digit, one carry         & $2 + 2$ & 1 & 5\,000 \\
7  & 2-digit + 2-digit, two carries       & $2 + 2$ & 2 & 5\,500 \\
8  & 2-digit + 3-digit, no carry          & $2 + 3$ & 0 & 5\,500 \\
9  & 2-digit + 3-digit, one carry         & $2 + 3$ & 1 & 6\,000 \\
10 & 3-digit + 3-digit, no carry          & $3 + 3$ & 0 & 6\,000 \\
\bottomrule
\end{longtable}

\paragraph{Notes.}
Carry count specifies exactly how many column additions overflow into the next column.
Level 1 guarantees single-digit sums (e.g.\ $3 + 4 = 7$); Level 5 requires handling two-digit
arithmetic with no carrying; Levels 6–10 extend to two- and three-digit operands with
increasing carry requirements.

% ──────────────────────────────────────────────
\section{Subtraction (SUB)}

\begin{longtable}{@{}clccr@{}}
\toprule
\textbf{Level} & \textbf{Description} & \textbf{Operands} & \textbf{Borrows} & \textbf{Cutoff (ms)} \\
\midrule
\endhead
1  & 1-digit $-$ 1-digit, non-negative          & $1 - 1$ & 0 & 4\,000 \\
2  & 2-digit $-$ 1-digit, no borrow             & $2 - 1$ & 0 & 4\,500 \\
3  & 2-digit $-$ 1-digit, with borrow           & $2 - 1$ & 1 & 4\,500 \\
4  & 2-digit $-$ 2-digit, no borrow             & $2 - 2$ & 0 & 5\,000 \\
5  & 2-digit $-$ 2-digit, with borrow           & $2 - 2$ & 1 & 5\,000 \\
6  & 3-digit $-$ 1-digit, no borrow             & $3 - 1$ & 0 & 5\,000 \\
7  & 3-digit $-$ 1-digit, with borrow           & $3 - 1$ & 1 & 5\,500 \\
8  & 3-digit $-$ 2-digit, no borrow             & $3 - 2$ & 0 & 5\,500 \\
9  & 3-digit $-$ 2-digit, with borrow           & $3 - 2$ & 1 & 6\,000 \\
10 & 3-digit $-$ 3-digit, no borrow             & $3 - 3$ & 0 & 6\,000 \\
\bottomrule
\end{longtable}

\paragraph{Notes.}
Negative results are never generated (\texttt{allowNegative: false} for all levels).
Level 1 constrains the generator so that the right operand is always $\leq$ the left operand,
avoiding borrows by construction.  Levels 6–10 introduce three-digit minuends, progressing
from simple subtraction up to full three-digit minus three-digit.

% ──────────────────────────────────────────────
\section{Multiplication (MUL)}

\begin{longtable}{@{}clccr@{}}
\toprule
\textbf{Level} & \textbf{Description} & \textbf{Operands} & \textbf{Carries} & \textbf{Cutoff (ms)} \\
\midrule
\endhead
1  & 1-digit $\times$ 1-digit, no carry          & $1 \times 1$ & 0 & 4\,000 \\
2  & 1-digit $\times$ 1-digit, with carry        & $1 \times 1$ & 1 & 4\,000 \\
3  & 1-digit $\times$ 2-digit, no carry          & $1 \times 2$ & 0 & 4\,500 \\
4  & 1-digit $\times$ 2-digit, one carry         & $1 \times 2$ & 1 & 4\,500 \\
5  & 1-digit $\times$ 2-digit, two carries       & $1 \times 2$ & 2 & 5\,000 \\
6  & 1-digit $\times$ 3-digit, no carry          & $1 \times 3$ & 0 & 5\,000 \\
7  & 1-digit $\times$ 3-digit, one carry         & $1 \times 3$ & 1 & 5\,500 \\
8  & 1-digit $\times$ 3-digit, two carries       & $1 \times 3$ & 2 & 5\,500 \\
9  & 1-digit $\times$ 3-digit, three carries     & $1 \times 3$ & 3 & 6\,000 \\
10 & 2-digit $\times$ 2-digit, no carry          & $2 \times 2$ & 0 & 6\,000 \\
\bottomrule
\end{longtable}

\paragraph{Notes.}
For multiplication, a carry occurs when a partial product exceeds 9 and must be added to the
next column.  Levels 6–9 extend to a three-digit multiplicand, with the carry count rising
from zero to three.  Level 10 introduces two-digit multipliers for the first time.

% ──────────────────────────────────────────────
\section{Division (DIV)}

All division questions are \textbf{exact integer} divisions (no remainder), and the quotient is
always $\geq 2$ (except Level 1 where the minimum quotient is 1).

\begin{longtable}{@{}clccr@{}}
\toprule
\textbf{Level} & \textbf{Description} & \textbf{Dividend / Divisor} & \textbf{Min.\ Quotient} & \textbf{Cutoff (ms)} \\
\midrule
\endhead
1  & 1-digit $\div$ 1-digit & $1 \div 1$ & 1 & 4\,000 \\
2  & 2-digit $\div$ 1-digit & $2 \div 1$ & 2 & 4\,500 \\
3  & 2-digit $\div$ 2-digit & $2 \div 2$ & 2 & 5\,000 \\
4  & 3-digit $\div$ 1-digit & $3 \div 1$ & 2 & 5\,500 \\
5  & 3-digit $\div$ 2-digit & $3 \div 2$ & 2 & 6\,000 \\
6  & 3-digit $\div$ 3-digit & $3 \div 3$ & 2 & 6\,500 \\
7  & 4-digit $\div$ 1-digit & $4 \div 1$ & 2 & 6\,500 \\
8  & 4-digit $\div$ 2-digit & $4 \div 2$ & 2 & 7\,000 \\
9  & 4-digit $\div$ 3-digit & $4 \div 3$ & 2 & 7\,500 \\
10 & 4-digit $\div$ 4-digit & $4 \div 4$ & 2 & 8\,000 \\
\bottomrule
\end{longtable}

\paragraph{Notes.}
Questions are generated by first choosing a valid divisor and quotient, then computing the dividend
as their product.  This guarantees an integer result.  The minimum quotient constraint prevents
trivially easy questions such as $n \div n = 1$.  Levels 6–10 scale up to four-digit dividends
with increasingly large divisors.

% ──────────────────────────────────────────────
\section{Summary: Cutoff Times}

The tables below show the cutoff time (in milliseconds) for each level across all four domains.

\noindent\textbf{Levels 1--5}
\begin{center}
\begin{tabular}{@{}lrrrrr@{}}
\toprule
\textbf{Domain} & \textbf{L1} & \textbf{L2} & \textbf{L3} & \textbf{L4} & \textbf{L5} \\
\midrule
Addition        & 4\,000 & 4\,000 & 4\,500 & 4\,500 & 5\,000 \\
Subtraction     & 4\,000 & 4\,500 & 4\,500 & 5\,000 & 5\,000 \\
Multiplication  & 4\,000 & 4\,000 & 4\,500 & 4\,500 & 5\,000 \\
Division        & 4\,000 & 4\,500 & 5\,000 & 5\,500 & 6\,000 \\
\bottomrule
\end{tabular}
\end{center}

\noindent\textbf{Levels 6--10}
\begin{center}
\begin{tabular}{@{}lrrrrr@{}}
\toprule
\textbf{Domain} & \textbf{L6} & \textbf{L7} & \textbf{L8} & \textbf{L9} & \textbf{L10} \\
\midrule
Addition        & 5\,000 & 5\,500 & 5\,500 & 6\,000 & 6\,000 \\
Subtraction     & 5\,000 & 5\,500 & 5\,500 & 6\,000 & 6\,000 \\
Multiplication  & 5\,000 & 5\,500 & 5\,500 & 6\,000 & 6\,000 \\
Division        & 6\,500 & 6\,500 & 7\,000 & 7\,500 & 8\,000 \\
\bottomrule
\end{tabular}
\end{center}

Division carries the longest cutoff times at higher levels because reasoning through multi-digit
division requires more working-memory steps than the equivalent addition or multiplication.
Addition, Subtraction, and Multiplication share the same cutoff schedule across Levels 6--10.

\end{document}
